% Precise editorial derivations for ED26-BM-RAT2021-10-3.
\begin{annexitem}{annex:p18.r103.reflexivity}{thm:p18-r103-reflexivity}{Reflexivity with the attitude bridge exposed}
\begin{proof}
From \(I(S,f,b)\) obtain \(d\), \(W(S,b)\) and
\(B(S,d\leftrightarrow b)\). The second AC1 conjunct gives
\(B(S,b\leftrightarrow r)\). Classical propositional consequence
of the two believed biconditionals includes both
\(d\leftrightarrow r\) and \(b\to r\). The explicitly assumed
belief closure therefore gives \(B(S,d\leftrightarrow r)\) and
\(B(S,b\to r)\). Apply the stipulated desire bridge to
\(W(S,b)\) and \(B(S,b\to r)\), obtaining \(W(S,r)\).
Together with \(d\) and \(B(S,d\leftrightarrow r)\), this is
exactly \(I(S,f,r)\) by the definition of intention.

The bridge is not dispensable. Consider a single belief-accessible
valuation at which \(d,b,r\) are all true. All the needed beliefs
and their classical closure then hold. Let the desire relation hold
for target proposition \(b\) and not for \(r\). The definition
gives \(I(S,f,b)\), and the AC1 forward premises hold, but
\(I(S,f,r)\) fails. This example deliberately does not impose the
desire bridge; it demonstrates that propositional belief closure alone
cannot supply the missing attitude. No empirical interpretation of this
valuation is needed or claimed.
\end{proof}
\end{annexitem}

\begin{annexitem}{annex:p18.r103.success-truth}{thm:p18-r103-success-truth}{The separate truth conditions of recognition, uptake and correctness}
\begin{proof}
Treat \(c=CA(S,H,f,p)\), \(r=K'(H,c)\),
\(b=B'(H,p)\), \(s=B(S,p)\) and \(p\) as the propositions
or attitude reports being compared, with knowledge of \(c\) requiring
\(c\) to be true. Recognition is consistent with \(r,c\) true
and \(b,p\) false. Thus even factive recognition of the attempt
does not force uptake or content truth. This is also consistent with
the speaker's false expectation that recognition would yield uptake:
belief in an equivalence does not require actual equivalence.

Uptake is consistent with \(b,c,r\) true and \(p,s\) false.
The speaker can want and expect the false belief and openly communicate
that content, without sincerity. Hence attainment of the informative
aim does not imply truth. Alternatively set \(s\) true while \(p\)
is false to represent a sincere but erroneous speaker; sincerity alone
does not imply correctness.

If \(c\), \(c\to s\) and \(s\to p\) hold, two applications
of modus ponens give \(s\) and \(p\). These premises contain no
condition forcing \(b\): a valuation with \(c,s,p,r\) true and
\(b\) false satisfies them. Consequently sincerity and correctness
suffice for truth of an asserted content but not for the hearer's belief.
These countervaluations concern the indicated implications, not the
existence of a complete model for the source's infinitary \(I^*\)
notation or the behavior of an observed population.
\end{proof}
\end{annexitem}

\begin{annexitem}{annex:p18.r103.common-meaning}{thm:p18-r103-common-meaning}{Common knowledge of a regularity is itself commonly known}
\begin{proof}
Let \(R_P\) be the reflexive-transitive closure of the union of
the epistemic accessibility relations of members of \(P\), as in
the cited handbook semantics. Suppose \(C_Pq\) holds at world
\(w\). Then \(q\) holds at every \(v\) with \(wR_Pv\).
Fix any such \(v\). If \(vR_Pu\), transitivity gives
\(wR_Pu\), so \(q\) holds at \(u\). Hence \(C_Pq\)
holds at \(v\). Since this is true for every \(R_P\)-reachable
\(v\), \(C_PC_Pq\) holds at \(w\). This proves
\(M\to C_PM\) for \(M=C_Pq\).

Reflexivity gives \(wR_Pw\), so \(q\) itself holds at \(w\),
proving \(M\to q\). Neither argument depends on the truth of
\(p\) or on the hearer's later belief: \(q\) is the regularity
of content-bearing attempts, not the sincerity or correctness condition.
The proof assumes a well-defined \(q\) and the declared group
semantics. It does not manufacture an interpretation of an arbitrary
intention hierarchy or derive common knowledge from merely one observed
regularity, or from a fixed finite nesting of individual knowledge.
\end{proof}
\end{annexitem}

\begin{annexitem}{annex:p18.r103.perfect-signal}{thm:p18-r103-perfect-signal}{The ideal signaling implication and the role of factivity}
\begin{proof}
From \(b\leftrightarrow d\) and \(d\leftrightarrow p\),
transitivity of material equivalence yields \(b\leftrightarrow p\).
If \(d\) holds, the second equivalence gives \(p\). The hypothesis
that the speaker knows whether \(p\) is
\(K(S,p)\lor K(S,\neg p)\). Factivity excludes the second
disjunct when \(p\) is true. Therefore \(K(S,p)\) holds, and
the stipulated knowledge-to-belief condition gives \(B(S,p)\).
The earlier conclusion \(p\) also supplies correctness of this belief.

If the joint strategy is common knowledge in the declared reflexive
semantics, it holds at the actual world, so the actual equivalences
used above follow. Replacing common knowledge by nonfactive common
belief invalidates this step. For example, take an actual world where
\(d,b\) are true and \(p\) false, and let all agents' belief
relations point only to one other world where \(d,b,p\) are true.
At both worlds the belief relations point to that other world. The
strategy equivalences are believed at every finite mutual-belief level,
but \(d\leftrightarrow p\) is false at the actual world. Thus common
belief in the strategy does not force actual truth transmission.
These relations are serial, transitive and Euclidean, so the example
does not depend on inconsistent or empty belief alternatives.

Finally, if \(d\) is false the strategy yields \(\neg b\), not
\(B'(H,\neg p)\). A belief set need not contain either of the
two attitudes. The proof therefore asserts exactly the positive-content
tracking condition stated in the source, not a stronger complete
two-answer reporting protocol.
\end{proof}
\end{annexitem}
